\section{Methodology}
\label{sec:methodology}

\subsection{Data}
\label{sec:data}
Three public MSI collections were used (\cref{tab:data}). They were chosen
because they are the collections on which S3PL was published, so both
baselines could be reproduced under their own protocols and then compared with
our methods under a common one.

\begin{table}[t]
\caption{Datasets. Sections are the units of analysis; the statistical unit for
inference is noted for each collection.}
\label{tab:data}
\centering\small
\begin{tabular}{@{}lccc@{}}
\toprule
 & GBM & CAC & RCC \\
\midrule
Sections & 8 & 8 & 8 \\
Independent units & 4 PDX lines & 1 biopsy & 8 subjects \\
$m/z$ bins & 85{,}062 & 1{,}481 & 10{,}200 \\
Pixels per section & 2{,}071--4{,}524 & 3{,}243--4{,}761 & 972--3{,}564 \\
Mask classes (GMM $K$) & 2 & 3 & 3 \\
\bottomrule
\end{tabular}
\end{table}

As planned in the proposal, the GBM data are the massNet PDX glioblastoma
dataset (PR001292, 9.4\,T FT-ICR)~\citep{abdelmoula2022massnet} and the CAC data
are the MetaboLights colorectal adenocarcinoma dataset
(MTBLS415)\checkme{acquisition paper for MTBLS415}. The expert masks come from
the S3PL benchmark release~\citep{weigand2026spatial}. RCC from \citet{dill2010rcc} via Cardinal~\citep{bemis2015cardinal}.
\textbf{GBM} contains eight sections from four patient-derived xenograft (PDX)
lines, two sections per line. Sections from the same line are not
independent, so GBM results are also summarised per PDX line.
\textbf{CAC} contains eight sections from a single biopsy, so its comparisons
are within one specimen. \textbf{RCC} contains eight runs from eight subjects
and has three classes: cancer, normal and a large unannotated region (45--78\%
of pixels). The dataset's own workflow treats the unannotated region as
background. We kept it as a third class in the primary analysis and report a
cancer/normal-only sensitivity analysis (\cref{sec:res-rcc}).
Measured-coordinate and mask alignment was audited for every section
(\cref{app:evidence}).

Four sections (GBM108\_positive, GBM22\_2, 160TopL and 40TopL) were used for
development and tuning. The method was frozen on 16 September, and these
sections are excluded from confirmation claims wherever the corresponding rule
specified it. RCC was not used during development and serves as a held-out
confirmation dataset.

\subsection{Pipeline}
\label{sec:pipeline}
For each section, a model is fitted to that section's spectra without labels,
and the selected peaks are then scored against that section's expert masks. The
design is therefore unsupervised and transductive. The expert masks never
enter model fitting or ranking. They do, however, set two configuration values:
the number of GMM components, and the peak budget. Both are disclosed below.

\textbf{VAE.} Spectra are TIC-normalised. Following msiPL, a fully connected
VAE with one hidden layer of 512 units and a 5-dimensional latent space is
trained for 100 epochs with Adam (learning rate $10^{-3}$, batch size 128). The
loss is msiPL's categorical cross-entropy reconstruction loss plus a KL term
($\beta = 1$)~\citep{abdelmoula2021peak,kingma2014autoencoding}. Two core arms
are compared:
\begin{itemize}
  \item \emph{centre-only}, which encodes one spectrum;
  \item \emph{uniform}, which concatenates the spectrum with the mean of its up
        to eight measured $3\times3$ neighbours.
\end{itemize}
Both reconstruct the centre spectrum only. For GBM ($D=85{,}062$) the two arms
have 87.2M and 130.8M parameters. The further variants are:
\begin{itemize}
  \item learned single-head attention over the neighbours, which adds about
        0.5\% parameters~\citep{vaswani2017attention};
  \item $5\times5$, $7\times7$ and $9\times9$ uniform windows;
  \item zeroed and nonlocally shuffled neighbours;
  \item centre masking, in which the centre is hidden for half of the training
        samples so that the neighbours become necessary.
\end{itemize}

\textbf{Segmentation.} A GMM with full covariance is fitted to the standardised
latent means. Its number of components $K$ equals the number of annotated
classes (\cref{tab:data}). This is a class-informed choice: BIC did not
support it (\cref{sec:res-encoder}), so the pipeline is label-free in fitting
and ranking but not fully autonomous.

\textbf{Rankings.} Every ranking returns an ordering of $m/z$ bins, and the top
$K_{\text{bin}}$ bins are scored.
\begin{itemize}
  \item \emph{Integrated Gradients} integrates the gradient of the assigned GMM
        component posterior with respect to the input spectrum. The path runs
        from the section's mean spectrum to each of 12 sampled pixels per
        component, using 64 trapezoidal steps over both the centre and the
        context inputs. Bins are ranked by mean absolute attribution, balanced
        across components~\citep{sundararajan2017axiomatic}.
  \item \emph{Posterior correlation} ranks bins by the absolute Pearson
        correlation between each bin's ion image and each component's posterior
        map, balanced across components in the same way. It uses the same GMM
        as IG and no gradients. This is the segment-then-correlate
        approach~\citep{alexandrov2011efficient}.
  \item \emph{Model-free rankings:} mean intensity, variance, and global
        Moran's~$I$ over measured Moore-neighbour edges~\citep{moran1950notes}.
  \item \emph{References:} a random ranking, and a supervised oracle that ranks
        by correlation with the expert masks. The oracle is a ceiling only and
        cannot be deployed.
\end{itemize}

\textbf{Controls on the representation and segmentation.}
\begin{itemize}
  \item \emph{Test~V} replaces the VAE latent space with a 5-component PCA
        while keeping the GMM and posterior correlation, which isolates the
        nonlinear encoder.
  \item \emph{Test~S} smooths the latent means with structure-adaptive
        neighbourhood weights before the GMM~\citep{alexandrov2011efficient},
        which places spatial information at the segmentation step instead of
        the encoder input.
\end{itemize}

\textbf{Baselines.} Legacy msiPL was reproduced from its released code and
tuned to its published peak-count procedure~\citep{abdelmoula2021peak}. S3PL
was run from the public repository with the released settings: 10 epochs, Adam
(learning rate $10^{-2}$), batch size 16, and patch size and peaks per patch of $p=3$, $z=256$ (GBM),
$p=9$, $z=256$ (CAC) and $p=5$, $z=1024$ (RCC)~\citep{weigand2026spatial}. The released code normalises each
patch by its spatial maximum. The paper describes TIC normalisation, so both
were tested (\cref{sec:res-s3pl}).

\subsection{Evaluation}
\label{sec:evaluation}
\textbf{mSCF1.} For each class mask and each threshold
$t\in\{0.3,0.4,0.5,0.6\}$, a reference set of bins is formed from the bins
whose ion images correlate with the mask. Following the released S3PL
evaluator, this uses the threshold-derived rank cutoff in each class's ranked
correlation list. The F1 score between the selected bins and the union of the
class reference sets is computed at each threshold, and mSCF1 is the mean over
the four thresholds~\citep{weigand2026spatial}. mSCF1 measures spatial
agreement with expert annotation, not chemical identity.

\textbf{Matched budgets.} All methods compared in fair scoring select the same
number of bins, $K_{\text{bin}}$, per section. For GBM and CAC this is the
tuned legacy msiPL peak count. For RCC it is the released S3PL target, which
approximately equals the number of bins correlating with the masks at
$t=0.5$. Budgets of $0.5\times$ to $2\times K_{\text{bin}}$ were checked as a
sensitivity analysis.

\textbf{Peak-level scoring (P1).} Neighbouring bins of one molecular peak
would otherwise each count as a true positive. P1 therefore collapses bins
within one basin of the section's mean spectrum into a single peak group.
P1 partitions were audited and approved before any score was computed.
On CAC and RCC, P1 groups are singletons or near-singletons, so P1 acts there
as a second budget check rather than independent evidence.

\subsection{Experimental design}
\label{sec:design}
Every comparison in \cref{sec:results} has a decision rule that was written
down before its outputs existed. Rules, deviations and execution logs were
kept in an append-only protocol (\cref{app:evidence}).
\begin{itemize}
  \item A difference ``helps'' if its mean is at least $+0.02$ mSCF1 and it is
        positive in at least 6 of 8 sections.
  \item Two methods are ``equivalent'' if the mean difference lies within
        $\pm0.02$.
  \item Anything else is ``mixed'', unless sections are missing
        (``incomplete'').
\end{itemize}
These practical labels are quoted exactly in the results.
Training seeds: all arms were trained with seed~1, and the two core arms were
retrained with seeds~2 and~3 on all 24 sections (Test~R). Collection-level
seed results are reported as mean $\pm$ SD across seeds of the seed-level
collection means. Per-section seed variability is reported as the median
(range) of per-section SDs.

\subsection{Statistical analysis}
\label{sec:stats}
The practical labels are accompanied by paired non-parametric statistics,
computed over sections:
\begin{itemize}
  \item the exact two-sided Wilcoxon signed-rank test~\citep{wilcoxon1945individual};
  \item the Hodges--Lehmann median paired difference with its exact 95\%
        confidence interval~\citep{hodges1963estimates};
  \item the matched-pairs rank-biserial correlation $r$;
  \item the sign count.
\end{itemize}
Where three seeds exist, seeds are averaged within each section first, so the
unit is the section and not the run. A Mann--Whitney test was not used for
these comparisons, because each pair of methods is scored on the same sections
and the data are therefore paired.

The 38 rule-labelled comparisons were grouped into five families by research
question: ranking (6 tests), versus legacy (4), encoder context and variants
(17), segmentation (9), and S3PL (2). Holm correction~\citep{holm1979simple} was
applied within each family. With $n=8$ the smallest attainable exact p-value is
$2/2^8=0.0078$, so Holm significance is attainable only in families of at most
six tests. Effect sizes and intervals therefore carry the main weight. A
comparison is called ``statistically supported'' only if its Holm-adjusted
$p<0.05$ \emph{and} its practical label points the same way. GBM is also
summarised per PDX line ($n=4$, smallest attainable $p=0.125$), outside the
Holm families. Mixed-effects models, with score as the outcome, method as a
fixed effect and random intercepts for section and seed, are reported
descriptively.

\subsection{Computational considerations}
\label{sec:compute}
All training ran on the University's MSCluster under Slurm, on RTX 3060 and
RTX 3090 GPUs. Ranking and scoring ran on CPU nodes.
\begin{itemize}
  \item \textbf{Caching.} The original data loader re-read each batch from HDF5
        on shared storage, and a 100-epoch GBM run took 11.3\,h. Profiling
        showed that data preparation took 24.6\,s per batch, while the forward
        and backward passes took about 0.02\,s. Reading each section into memory
        once (about 13\,s) reduced the run to 23\,min, a $30\times$ speed-up,
        with identical loss histories and final weights. The likely mechanism is
        repeated small reads over network storage, which was not separately
        measured.
  \item \textbf{Bit-identity gates.} Before reusing saved runs, same-seed
        reruns were checked for exact equality of losses, scores and selected
        peaks.
  \item \textbf{Hardware and exact equality.} Exact equality of rankings
        depends on the CPU model, because floating-point summation order
        differs between processors. Timing runs were therefore matched to the
        processor that produced each saved ranking.
\end{itemize}
Runtimes are reported in \cref{sec:res-runtime}.
